\documentclass[10pt,a4paper]{amsart}
\usepackage[margin=19mm]{geometry}
\usepackage{amsmath,amssymb,mathtools}
\usepackage[scr=rsfs,cal=euler]{mathalfa}
\usepackage{xfrac}
\usepackage[unicode,hidelinks,bookmarks=false]{hyperref}
\newcommand{\ie}{i.e.\ }
\newcommand{\cf}{cf.\ }
\newcommand{\resp}{respectively\ }
\newcommand{\Subsection}{\subsection}
\providecommand{\nobreakdash}{\nobreak\hbox{-}}
\setlength{\emergencystretch}{2em}
\multlinegap0pt
\usepackage{amssymb}
\usepackage{enumerate}
\newtheorem{thm}{Theorem}
 \def\Ext{\operatorname{Ext}}
\def\Hom{\operatorname{Hom}}
\newcommand{\NN}{\mathbb{N}}
\newcommand{\fa}{\mathfrak{a}}
\def\grade{\operatorname{grade}}
\newcommand{\m}{\mathfrak{m}}
\newcommand{\p}{\mathfrak{p}}
\def\var{\operatorname{Var}}
\title{Bass numbers of local cohomology modules at the first and last non-vanishing levels}
\author{M. Jahangiri, R. Ahangari Maleki}
\date{Source: arXiv:2603.18724v1 (CC BY 4.0)}
\begin{document}
\maketitle

\noindent\emph{Source and licence.} Bass numbers of local cohomology modules at the first and last non-vanishing levels. Version arXiv:2603.18724v1 (CC BY 4.0), distributed by arXiv under the Creative Commons Attribution 4.0 International licence, http://creativecommons.org/licenses/by/4.0/, as recorded in the arXiv licence metadata for this exact version and shown on the abstract page https://arxiv.org/abs/2603.18724v1. Full licence notice: \texttt{licenses/arxiv-2603.18724-CC-BY-4.0.txt}.\par\smallskip

\noindent\textit{Editorial scope.} Counted unit: the article's thm grade (source0.tex lines 263--271). The article's complete original proof of it is delivered with it (source0.tex lines 273--318, one \texttt{proof} environment of 46 lines). No mathematics is written, repaired or re-derived: the statement and the proof are the article's own lines, in the article's own notation, and no retained line is substituted or re-flowed.  Retained uncounted as the source-local context the proof needs: lines 227--229.  Nothing was omitted from any retained span, so the package carries no omission record.  The retained lines cite 1 works, whose entries are transcribed verbatim from the article's own bibliography source in the e-print and delivered with short editorial labels recorded in the item's reference mapping.  No retained line contains a figure environment, an image-inclusion command or drawing code, so no image is delivered and the package declares no asset dependency.  Editorial formatting only, in addition: an AMS \texttt{amsart} preamble that restates the article's own declarations and macros used by the retained lines, namely the environments \texttt{thm} on separate counters, together with the standard packages those lines need.  The retained environments print in this document as Theorem 1 in order of appearance, whereas the article prints the same items with the numbering of its own full document.  The complete native source of the article as distributed in the arXiv e-print for this exact version is bundled unchanged as \texttt{sources/196714/source0.tex}.
\begin{equation}\label{1}
     E_2^{i, j}= \Ext^i_{R }(\tfrac{R}{\m}, H^j_{\fa}(N))\underset{i}{\Rightarrow}\Ext^{i+ j}_{R}(\tfrac{R}{\m}, N).
\end{equation}

\begin{thm}\label{grade}
Let $\p\in \var(\fa)$ and   $g\leq \grade_{\fa}(M)$. Then the following statements hold.
\begin{enumerate}
    \item $\mu^0(\p, H^g_{\fa}(M))= \mu^g(\p, M)< \infty$.
    \item $\mu^1(\p, H^g_{\fa}(M))\leq \mu^{g+ 1}(\p, M)$ and if $\mu^0(\p, H^{g+ 1}_{\fa}(M))= 0$, then the equality holds.
    \item If $\mu^2(\p, H^g_{\fa}(M))= 0$, then $\mu^0(\p, H^{g+ 1}_{\fa}(M)) \leq \mu^{g+ 1}(\p, M)< \infty$.
    \item If $\mu^3(\p, H^g_{\fa}(M))= 0$, then $\mu^1(\p, H^{g+ 1}_{\fa}(M)) \leq \mu^{g+ 2}(\p, M)< \infty$.
\end{enumerate}
    \end{thm}

\begin{proof}
First of all note that, $\grade_{\fa}(M)\leq \grade_{\fa R_{\p}}(M_{\p})$ and since $\p\in \var(\fa)$, by  the flat base change theorem (\cite[Theorem 4. 3. 2]{bsh}), after replacing $R$ with $R_{\p}$ we may assume that $(R, \m)$ is local and $\p= \m$. Now, consider the convergence of spectral sequences \ref{1},
\[ E_2^{i, j}= \Ext^i_{R }(\tfrac{R}{\m}, H^j_{\fa}(M))\underset{i}{\Rightarrow}\Ext^{i+ j}_{R}(\tfrac{R}{\m}, M).\]
 \begin{enumerate}
     \item Since $g\leq \grade_{\fa}(M)$, by \cite[Theorem 6. 2. 7]{bsh},
     \begin{equation}\label{zero}
       0= \Ext^i_{R }(\tfrac{R}{\m}, H^j_{\fa}(M))= E_2^{i, j}=  E_{\infty}^{i, j} \,\,\,\,\text{for all}\,\,i\in \NN_0\,\,\text{and all}\,\,j< g.
     \end{equation}
     Also, there exists a filtration 
 \[0\subseteq \Phi_{g} \subseteq\cdots\subseteq \Phi_{   1} \subseteq\Phi_{0}= \Ext^{g}_{R}(\tfrac{R}{\m}, M),\]
 of submodules of $\Ext^{g}_{R}(\tfrac{R}{\m}, M)$ such that 
 \begin{equation*}
     E_{\infty}^{i, j}\cong \tfrac{\Phi_{i}}{\Phi_{i+ 1}},\,\,\text{for all}\,\,i, j\in \NN_0\,\,\text{with}\,\,i+ j= g.
 \end{equation*}
 Now, \ref{zero} implies that 
 \[\Hom_{R }(\tfrac{R}{\m}, H^g_{\fa}(M))=  E_2^{0, g}= E_{\infty}^{0, g}\cong  \Ext^{g}_{R}(\tfrac{R}{\m}, M).\]
 \item Using \ref{zero} we have 
 \[\Ext^1_{R }(\tfrac{R}{\m}, H^g_{\fa}(M))= E_2^{1, g}=  E_{\infty}^{1, g}.\]
 Furthermore, there exists a filtration 
 \begin{equation}\label{f1}
     0\subseteq \Phi_{g+ 1} \subseteq\cdots\subseteq \Phi_{   2} \subseteq\Phi_{   1} \subseteq\Phi_{0}= \Ext^{g+ 1}_{R}(\tfrac{R}{\m}, M),
 \end{equation}
 of submodules of $\Ext^{g+ 1}_{R}(\tfrac{R}{\m}, M)$ such that 
 \begin{equation*}
     E_{\infty}^{i, j}\cong \tfrac{\Phi_{i}}{\Phi_{i+ 1}},\,\,\text{for all}\,\,i, j\in \NN_0\,\,\text{with}\,\,i+ j= g+ 1.
 \end{equation*}
 \ref{zero} implies that the filtration \ref{f1} is actually of the form
 \[ 0= \Phi_{g+ 1} =\cdots= \Phi_{   2} \subseteq\Phi_{   1} \subseteq\Phi_{0}= \Ext^{g+ 1}_{R}(\tfrac{R}{\m}, M).\]
 Therefore,
\begin{equation}\label{2}
    \Ext^1_{R }(\tfrac{R}{\m}, H^g_{\fa}(M))\cong \Phi_{   1},
\end{equation}
 is a submodule of $\Ext^{g+ 1}_{R}(\tfrac{R}{\m}, M)$. If $\mu^0(\p, H^{g+ 1}_{\fa}(M))= 0$ then 
 \[0= E_{\infty}^{0, g+ 1}= \tfrac{\Ext^{g+ 1}_{R}(\tfrac{R}{\m}, M)}{\Phi_{  1}}.\]
 Hence, using \ref{2} we have 
 \[ \Ext^1_{R }(\tfrac{R}{\m}, H^g_{\fa}(M))\cong \Ext^{g+ 1}_{R}(\tfrac{R}{\m}, M),\]
 and the result follows.
 
 \item Assume that $\mu^2(\p, H^g_{\fa}(M))= 0$. Then $E_2^{2, g}= 0$ and by \ref{zero} we deduce that 
 \[\Hom_{R }(\tfrac{R}{\m}, H^{g+ 1}_{\fa}(M))= E_2^{0, g+ 1}= E_{\infty}^{0, g+ 1}.\]
 Using the concept of convergence of spectral sequences, as used in parts (1) and (2),  the latter module is a subquotient of $\Ext^{g+ 2}_{R}(\tfrac{R}{\m}, M)$. This proves the claim.
 \item Let $\mu^3(\p, H^g_{\fa}(M))= 0$. Then $E_2^{3, g}= 0$ and by \ref{zero} we have 
 \[\Ext^1_{R }(\tfrac{R}{\m}, H^{g+ 1}_{\fa}(M))= E_2^{1, g+ 1}= E_{\infty}^{1, g+ 1}.\] 
 Therefore, using similar argument as in the above parts,  $\Ext^1_{R }(\tfrac{R}{\m}, H^{g+ 1}_{\fa}(M))$ is a subquotient of $\Ext^{g+ 2}_{R}(\tfrac{R}{\m}, M)$. In other words, $$\mu^1(\m, H^{g+ 1}_{\fa}(M)) \leq \mu^{g+ 2}(\m, M)< \infty.$$
 \end{enumerate}
\end{proof}

\begin{thebibliography}{99}
\bibitem[bsh]{bsh} M.\ Brodmann, R.Y.\ Sharp, {\em Local Cohomology: An Algebric Introduction with Geometric Applications}, 2nd ed. Cambridge University Press, (2012).
\end{thebibliography}
\end{document}
